% 1-6-scope.tex
%  Budget: ~1.3pp (detailed, per request). Own design -- only lago is cited.
%  KEY: exclusions (fuel/weather/cross-border) are a CONSIDERED decision that
%  STRENGTHENS the benchmark claim (WRITING_HANDOFF.md 8), NOT an oversight.
%  No invented numbers: OOD window qualitative; benchmark span structural.
%  Consistent terms: روز-پیش, بار پایه‌ی روزانه, «آرام»/«پرتنش».

\section{محدوده و مفروضات کلی پروژه}

برای آن‌که این پژوهش کانونی، و نتایج آن تکرارپذیر و منصفانه باقی بماند، دامنه‌ی
آن به‌عمد محدود شده است. مطالعه بر بازار برق آلمان (\lr{EPEX-DE}) و تنها بر افق
پیش‌بینیِ روز-پیش (\lr{day-ahead}) متمرکز است و بازارهای درون‌روز و تعادلی را در
بر نمی‌گیرد. کانونِ کار، پیش‌بینیِ نقطه‌ای است و نه پیش‌بینیِ بازه‌ای یا
احتمالاتی؛ و دو هدف دنبال می‌شود: بردار قیمت‌های ساعتیِ روز آینده و بار پایه‌ی
روزانه.

داده‌های مورد استفاده در دو دسته جای می‌گیرند. بخش اصلیِ مطالعه بر مجموعه‌داده‌ی
بنچمارکِ عمومیِ لاگو و همکاران \cite{lago_forecasting_2021} برای این بازار استوار
است که مقایسه‌ی مستقیم با نتایج منتشرشده را ممکن می‌سازد و یک بازه‌ی چندساله را
با دوره‌ی آزمونِ کنارگذاشته‌ی دوساله در بر می‌گیرد. افزون بر آن، برای آزمون
پایداریِ برون‌توزیعی، از حدود شش ماه داده‌ی زنده‌ی سال \lr{2026} از منبع
\lr{Energy-Charts} (مؤسسه‌ی فراونهوفر \lr{ISE}) استفاده می‌شود؛ این داده تنها برای
ارزیابیِ برون‌توزیعی به‌کار می‌رود و هیچ نقشی در آموزش مدل‌ها ندارد.

از نظر مدل‌ها، مجموعه‌ای منتخب و ثابت بررسی می‌شود که طیفی از یک مدل ساده‌ی مبنا،
مدل‌های آماری، مدل‌های یادگیری ماشین و یادگیری عمیق، و روش‌های ترکیبی را در بر
می‌گیرد. هدف، مقایسه‌ی نظام‌مندِ این خانواده‌ها در یک چارچوب یکسان است، نه مرورِ
جامعِ همه‌ی معماری‌های ممکن.

در انتخاب ورودی‌ها نیز محدودیتی آگاهانه اعمال شده است. تنها متغیرهای برون‌زایی
به‌کار می‌روند که در چارچوب پروتکل بنچمارک در دسترس‌اند \lr{—} یعنی پیش‌بینیِ
روز-پیشِ بار مصرفی و پیش‌بینیِ روز-پیشِ تولید (عمدتاً از منابع تجدیدپذیر) \lr{—}
در کنار مقادیر باوقفه‌ی خودِ قیمت. در مقابل، متغیرهایی مانند قیمت سوخت (گاز و
زغال‌سنگ)، متغیرهای جوی و جریان‌های تبادلیِ مرزی به‌عمد از دامنه‌ی پژوهش کنار
گذاشته شده‌اند. این کنارگذاری یک تصمیمِ سنجیده است و نه یک غفلت: افزودنِ
داده‌های برون‌زای فراتر از پروتکل، مقایسه‌ی منصفانه با نتایج منتشرشده را ناممکن
می‌کرد، و افزون بر آن برخی از این سری‌ها \lr{—} به‌ویژه قیمت سوخت \lr{—} به‌صورت
باز و تمیز در دسترس نیستند. بدین‌ترتیب، این محدودسازی به‌جای آن‌که ضعف به‌شمار
آید، اعتبارِ ادعای مقایسه‌ایِ پژوهش را تقویت می‌کند.

این پژوهش بر چند فرضِ بنیادیِ زیر استوار است:

\begin{enumerate}
\item متغیرهای برون‌زای مورد استفاده، پیش‌بینی‌های روز-پیشِ در دسترس پیش از
بسته‌شدن بازارند و نه مقادیر محقق‌شده؛ این فرض با شرایط واقعیِ بهره‌برداری و با
پروتکل بنچمارک سازگار است.

\item قیمت‌های منفی بخشی حقیقی از رفتار این بازارند و حذف یا برش داده نمی‌شوند؛ در
نتیجه، معیارهای خطایی که در حضور مقادیر صفر یا منفی نامعتبر می‌شوند (مانند
\lr{MAPE}) کنار گذاشته می‌شوند و به‌جای آن‌ها معیارهای مبتنی بر خطای مطلق به‌کار
می‌روند.

\item مجموعه‌داده‌ی بنچمارک به‌همان شکلِ منتشرشده و به‌مثابه‌ی مرجعِ متعارف
پذیرفته می‌شود تا تکرارپذیریِ مقایسه حفظ شود.

\item رفتار بازار میان رژیم‌های «آرام» و «پرتنش» ناهمگن است؛ همین ناهمگنی، مبنای
رویکردِ گروهیِ رژیم‌محور در این پژوهش است.

\item در آزمون پایداریِ برون‌توزیعی، مدل‌ها پس از آموزش بر داده‌ی بنچمارک
«منجمد» می‌مانند و بازآموزی نمی‌شوند، تا هر تغییر در عملکرد، بازتابِ تغییرِ توزیعِ
داده باشد و نه سازگاریِ دوباره‌ی مدل با شرایط تازه.
\end{enumerate}
